\documentclass[11pt,a4paper]{article}

\usepackage{amsmath,amssymb,amsthm}
\usepackage{hyperref}
\usepackage{geometry}
\usepackage{doi}
\geometry{margin=2.5cm}

% ---- theorem environments ----
\newtheorem{theorem}{Theorem}[section]
\newtheorem{proposition}[theorem]{Proposition}
\newtheorem{lemma}[theorem]{Lemma}
\newtheorem{corollary}[theorem]{Corollary}
\newtheorem{definition}[theorem]{Definition}
\theoremstyle{remark}
\newtheorem{remark}[theorem]{Remark}
\newtheorem{openproblem}[theorem]{Open Problem}

% ---- notation shorthands ----
\newcommand{\bstar}{\beta^{*}}
\newcommand{\dpair}{\delta_{\mathrm{pair}}}
\newcommand{\dc}{\delta_{c}}
\newcommand{\sigmapair}{\sigma_{\mathrm{pair}}}
\newcommand{\sigmacan}{\sigma_{\mathrm{pair}}^{\mathrm{can}}}
\newcommand{\Dcan}{D^{\mathrm{can}}}
\newcommand{\rcq}{r(c,q)}
\newcommand{\Heis}{\mathrm{Heis}_{3}(\mathbb{Z}/q\mathbb{Z})}
\newcommand{\ellgam}{\ell_{\gamma}}


\title{Canonical Pair Observables in Weil Blocks:\\
The Residual Amplitude Factor and a Proposed Canonical
Normalisation with Open Tests\\[6pt]
{\large O19 --- Spectral Admissibility Sub-programme}}

\author{J.~Beau\\[4pt]
\small Independent Researcher}

\date{2026\\Version 1.3}


\hypersetup{
  pdftitle={Canonical Pair Observables in Weil Blocks: The Residual Amplitude Factor
    and a Proposed Canonical Normalisation with Open Tests},
  pdfauthor={Jérôme Beau}
}

\begin{document}
  \maketitle

  \begin{abstract}
    The O12--O18 sequence identifies the pair-level quantity
    $\sigmapair(n) = \sigma_{c}(n)\,\sigma_{q-c}(n)$ as the structurally motivated
    candidate observable on $\Heis$; a capacity exponent near
    $\dpair \approx 7.44$ was historically compared with the range
    $\dpair \in [7.4, 10.6]$, a value O16 no longer asserts.
    A residual amplitude factor $\rcq \in \{1,2,3,4\}$, depending on the Gram--Schmidt
    initialisation parameters $(b_{1}, b_{2})$, was reported in earlier work; those values are not
    reproduced, and on the real O12/O13 pipeline the exact conjugate-block identity of O16 forces
    $\rcq \equiv 1$ within O16's contract.
    A non-trivial factor is treated here as a hypothesis about a pre-saturation quantity.
    The present paper begins by making explicit the implicit normalisation factor
    $D(c, b_{1}, b_{2}, n)$ of the O12/O13 implementation, separating its
    dependence on the initial basis, the block instance, and the normalisation
    convention.
    We then formulate two independence tests for $\rcq$: variation under basis
    change at fixed block instance (Test~P2), and variation under block-instance
    change at rigidly fixed basis (Test~P3).
    Three non-trivial outcomes are considered, alongside a fourth (an invariant of $(c,q)$,
    including $\rcq \equiv 1$), which is the one that holds for the factor as defined here: $\rcq$ is a pure basis artefact (Case~A,
    eliminable by canonical choice), a discrete instance invariant (Case~B,
    requiring a canonical quotient), or a mixed quantity with both components
    (Case~C, requiring a decomposition $r = r_{\mathrm{base}}\,r_{\mathrm{inst}}$).
    For the residual factor as defined here, an asymptotic limit, the cumulative spans reach $q$
    on the finite Cayley graph, so $\rcq \equiv 1$ for every seed and block instance
    (Lemma~2.4): the fourth case holds and Cases~A--C cannot occur.
    They, and the tests, remain meaningful only for a residual factor defined on a pre-saturation
    window, which is not defined here and for which the tests are open.
    For each case we give the construction of a canonical observable $\sigmacan(n)$ that would
    follow and relate it heuristically to the metaplectic central-phase coherence $\ellgam(n)$
    of~O12; its independence from pipeline conventions is stated as an open problem, not proved.
    The exponent $\dpair$ is unaffected in all cases; the contribution of this
    paper concerns the amplitude structure, which is a prerequisite for the
    fibre-level study of the historical $\dpair \to \bstar$ prescription deferred to~O21.
  \end{abstract}

  \paragraph{Keywords.}
    spectral admissibility, non-injective projection, Weil representation,
    Gram–Schmidt process, canonical normalisation, pair observables, amplitude structure

    \clearpage
    \tableofcontents

    \clearpage

  \section{Introduction}
  \label{sec:intro}


  The O-series spectral admissibility sub-programme has progressively refined
  the observable on Heisenberg Cayley graphs
  $\mathrm{Cay}(G_{q}, S_{q})$ with $G_{q} = \Heis$ for prime~$q$.
  Papers O12--O15~\cite{Beau2026a16,Beau2026a17,Beau2026a18,Beau2026a19}
  established the exact block-level exponent
  $\hat{\delta}_{\mathrm{exact}} \approx 3.72$, structurally below the
  historical comparison range $\delta \in [7.4, 10.6]$.
  O16~\cite{Beau2026a20} proposed that the physically relevant
  observable lives at the level of conjugate pairs $\{c, q-c\}$,
  with $\dpair = 2\dc$ when the conjugate ratio is shell-independent; the historical value near $7.44$ is not asserted by O16.
  Within its own toy model, O17~\cite{Beau2026a21} shows that the raw Gram--Schmidt
  dynamics of blocks $c$ and $q-c$ coincide, so that any variation of the factor $\rcq$ would
  reside in the normalisation; none of this is asserted there for the O12/O13 pipeline, on
  which the equality of the conjugate exponents follows instead from O16's conjugate-block
  identity within O16's contract.
  O18~\cite{Beau2026a22} establishes Born--Infeld parity as a covariance of the substrate action's response family
  and states the identification of the parity orbit $\{c, q-c\}$ with a fibre of $\Pi$ as an open problem;
  the pair observable of O16--O17 therefore retains the status of a structurally motivated hypothesis.

  A residual amplitude factor $\rcq \in \{1, 2, 3, 4\}$ of the pair observable
  $\sigmapair(n)$, depending on the pipeline normalisation $D(c, b_{1}, b_{2}, n)$
  implicit in the O12/O13 implementation~\cite{Beau2026a16,Beau2026a17}, was reported in
  earlier versions of O16.
  Three statements must be kept apart.
  The integer values are a historical hypothesis whose source is not reproduced.
  O16~\cite{Beau2026a20} shows that, on the real O12/O13 pipeline, its exact conjugate-block
  identity forces $\rcq \equiv 1$; this holds within O16's contract and is not a statement
  about other observables or normalisations.
  O17~\cite{Beau2026a21} shows, within its own toy model only, that any variation of $\rcq$
  would reside entirely in the normalisation $D$.
  For its own asymptotic definition of $\rcq$, the present paper shows that $\rcq \equiv 1$
  (Lemma~\ref{lem:rank-seed} and the consequence stated after Remark~\ref{cor:D-asymptotic}); a
  non-trivial factor is treated as a hypothesis about a pre-saturation quantity.
  O18 Section~4.3 lists the construction of a canonical normalisation
  independent of $(b_{1}, b_{2})$ among its open directions, deferred
  to the present paper.

  The goal of O19 is not to assume in advance that $\rcq$ is eliminable.
  Instead, we formulate the tests that would determine its dependence structure, and give
  the canonical observable that each outcome would call for.
  The dependence structure of $\rcq$ is not assumed a priori; besides a fourth possibility
  ($\rcq$ an invariant of $(c,q)$, including $\rcq \equiv 1$), it may follow three distinct
  scenarios: a pure basis artefact (Case~A, eliminable by canonical choice of initial frame),
  a discrete instance invariant (Case~B, requiring a canonical quotient), or a mixed quantity
  admitting a decomposition $\rcq = r_{\mathrm{base}} \cdot r_{\mathrm{inst}}$ with both
  components (Case~C).
  All three cases are treated; for the factor of Definition~\ref{def:D} the classification is
  settled by Lemma~\ref{lem:rank-seed} (the fourth case holds), and for a pre-saturation factor it
  would be determined empirically by the tests of Sections~\ref{sec:testP2}--\ref{sec:testP3},
  which are posed here as open problems and not carried out.

  \paragraph{Scope.}
    O19 does not aim to modify the observable class or the exponent, but to state how the last
    layer of pipeline dependence at the amplitude level could be removed.
    It does not re-derive $\dpair$, does not address the open fibre-identification problem of O18, and does
    not address the historical $\dpair \to \bstar$ prescription, which is the subject of~O21.

  \paragraph{Organisation.}
    Section~\ref{sec:pipeline} makes explicit the normalisation factor
    $D(c, b_{1}, b_{2}, n)$.
    Section~\ref{sec:phase} relates its structure to the metaplectic central-phase
    coherence of the Weil representation.
    Sections~\ref{sec:testP2} and~\ref{sec:testP3} present the two independence
    tests.
    Section~\ref{sec:classification} lists the possible outcomes, Case~A, B, or~C.
    Section~\ref{sec:canonical} gives the case-dependent construction of the canonical observable.
    Section~\ref{sec:discussion} interprets the result.
    Section~\ref{sec:stability} verifies that $\dpair$ is unaffected.
    Section~\ref{sec:conclusion} concludes.

  \section{Explicit Form of the Pipeline Normalisation}
  \label{sec:pipeline}

  \subsection{Setup and notation}

    Let $q$ be an odd prime and $c \in (\mathbb{Z}/q\mathbb{Z})^{\times}$.
    The Weil representation $\rho_{c}$ of $\Heis$ acts on
    $\mathbb{C}^{q}$~\cite{Beau2026a20,Beau2026a21}.
    For a BFS shell $S_{n}$ of depth~$n$ in $\mathrm{Cay}(G_{q}, S_{q})$, the
    Gram--Schmidt span tracker builds an orthonormal basis of
    $\mathrm{span}\{\rho_{c}(g)v_{0} : g \in B_{n}\}$ iteratively, where
    $v_{0} \in \mathbb{C}^{q}$ is an initial seed vector and $B_{n}$ is the
    ball of radius~$n$.

    The O12/O13 implementation records at each shell $S_{n}$ the increment
    $\Delta\Sigma^{(c)}(n) = \dim\,\mathrm{span}(B_{n}) - \dim\,\mathrm{span}(B_{n-1})$,
    from which the cumulative span $\Sigma^{(c)}(n)$ and the observable
    $\sigma_{c}(n)$ are derived.
    Throughout this paper we maintain the following notational distinction:
    $\Sigma^{(c)}(n) \in \mathbb{Z}_{\geq 0}$ denotes the raw cumulative span (an integer,
    the dimension of the spanned subspace), while $\sigma_{c}(n)$ denotes the
    normalised observable derived from it (e.g., divided by a shell-size or
    dimension factor as specified in the O12/O13 pipeline~\cite{Beau2026a16}).
    The amplitude factor $\rcq$ originates in $\Sigma^{(c)}$ and propagates unchanged
    into $\sigma_{c}$.
    This distinction is essential: the exponent $\dc$, fitted on a pre-saturation window, is defined at the level
    of $\sigma_{c}(n)$, while the residual factor $\rcq$ originates at the level of
    $\Sigma^{(c)}(n)$.

    The factor $D(c, b_{1}, b_{2}, n)$ is implicit in this construction.
    Its explicit form can be written as follows.
    Let $\{e_{1}, \ldots, e_{k(n)}\}$ be the orthonormal basis produced by the
    Gram--Schmidt procedure on $\{\rho_{c}(g)\,v_{0} : g \in B_{n}\}$, where
    $v_{0} = v_{0}(b_{1}, b_{2})$ is the seed vector determined by the block
    instance and $k(n) = \Sigma^{(c)}_{(b_{1},b_{2})}(n)$.
    At each new shell element $g \in S_{n}$, the Gram--Schmidt step contributes
    \[
      \Delta_{g}
      \;=\;
      \left\|
        \rho_{c}(g)\,v_{0}
        \;-\;
        \sum_{j=1}^{k(n-1)} \langle \rho_{c}(g)\,v_{0},\, e_{j} \rangle\, e_{j}
      \right\|^{2},
    \]
    and $\Delta\Sigma^{(c)}(n) = \#\{g \in S_{n} : \Delta_{g} > 0\}$.

    \begin{lemma}[Phase invariance of $\Delta_{g}$]
      \label{lem:Deltag-phase}
      The residual norm $\Delta_{g}$ is invariant under multiplication of $v_{0}$ by a global
      phase $e^{i\theta}$.
      Under a unitary change of basis $v_{0} \mapsto U v_{0}$ with $U \in \mathrm{U}(q)$,
      $\Delta_{g}$ changes if and only if the projection of $\rho_{c}(g)\,Uv_{0}$ onto the
      current orthonormal frame differs from that of $\rho_{c}(g)\,v_{0}$.
      Consequently, any variation of $D(c, b_{1}, b_{2}, n)$ under Test~P2
      captures only the misalignment between $v_{0}$ and the orbit geometry of $\rho_{c}$,
      not an intrinsic property of the representation.
    \end{lemma}

    \noindent
    This lemma closes the link between $\Delta_{g}$ and Test~P2: a variation of $\rcq$
    under unitary changes of $v_{0}$ is evidence of a phase-alignment effect, not of a
    representation-theoretic invariant.
    The pipeline normalisation factor is then
    \begin{equation}
      \label{eq:D-explicit}
      D(c, b_{1}, b_{2}, n)
      \;=\;
      \frac{\Sigma^{(c)}_{(b_{1},b_{2})}(n)}{\Sigma^{(c)}_{\mathrm{ref}}(n)},
    \end{equation}
    where $\Sigma^{(c)}_{\mathrm{ref}}(n)$ is a reference span independent of
    $(b_{1}, b_{2})$, to be specified in Section~\ref{sec:canonical}.
    Equation~\eqref{eq:D-explicit} makes explicit what was implicit in O12/O13:
    $D$ encodes how the choice of seed $v_{0}$ aligns with the orbit structure
    of $\rho_{c}$ on $B_{n}$.

    \begin{definition}[Pipeline normalisation factor]
      \label{def:D}
      The \emph{pipeline normalisation factor} is $D(c, b_{1}, b_{2}, n)$ as given
      by equation~\eqref{eq:D-explicit}.
      The \emph{residual amplitude factor} is defined as the limit of the log-ratio:
      \begin{equation}
        \label{eq:r-def}
        \log \rcq
        \;:=\;
        \lim_{n\to\infty}
        \Bigl(
        \log \Sigma^{(c)}_{(b_{1},b_{2})}(n)
        \;-\;
        \log \Sigma^{(q-c)}_{(b_{1},b_{2})}(n)
        \Bigr),
      \end{equation}
      whenever it exists.
      Equivalently, $\rcq = \lim_{n\to\infty} D(c,b_{1},b_{2},n) / D(q-c,b_{1},b_{2},n)$ when the two
      reference spans have ratio tending to~$1$ (as they do when both reach~$q$).
      The convergence is in the ordinary sense (not limsup).
      By Lemma~\ref{lem:rank-seed} below, the limit exists and equals~$0$, so that $\rcq = 1$ for
      every block instance and seed.
      Independently, on O16's contract, where the conjugate ratio is taken on the exact block
      observable $\Sigma_n^{(c_1,c_2,c_3)}$ of the O12/O13 pipeline, O16's conjugate-block identity
      also gives~$1$.
    \end{definition}

    \begin{remark}
      \label{rem:D-layers-sep}
      Definition~\ref{def:D} separates two logically distinct objects: the shell-by-shell ratio $D$,
      which depends on $n$, and the asymptotic residual $\rcq$, which is $n$-independent by
      construction.
      Versions up to~1.2 left the dependence of the limit~\eqref{eq:r-def} on $(b_{1}, b_{2})$ to
      the tests of Sections~\ref{sec:testP2}--\ref{sec:testP3}; by Lemma~\ref{lem:rank-seed} the
      limit does not depend on them.
      The tests remain meaningful only for a residual factor defined on a pre-saturation window,
      which this paper does not define.
    \end{remark}

  \subsection{Closed-form expression for $D$ and its seed dependence}
    \label{subsec:D-closed}

    We derive a closed-form expression for $D(c, b_{1}, b_{2}, n)$ in terms of the
    Gram--Schmidt process, making the seed dependence explicit.

    Let $\mathcal{O}_{n} = \{\rho_{c}(g)\,v_{0} : g \in B_{n}\} \subset \mathbb{C}^{q}$
    be the orbit set at depth~$n$.
    The Gram--Schmidt procedure applied to $\mathcal{O}_{n}$ in a fixed enumeration order
    $g_{1}, g_{2}, \ldots$ produces an orthonormal sequence
    $e_{1}, e_{2}, \ldots, e_{k(n)}$ with $k(n) = \Sigma^{(c)}_{(b_{1},b_{2})}(n)$.
    The $j$-th basis vector is
    \[
      e_{j}
      \;=\;
      \frac{w_{j}}{\|w_{j}\|},
      \qquad
      w_{j} = \rho_{c}(g_{j})\,v_{0} - \sum_{i < j} \langle \rho_{c}(g_{j})\,v_{0},\, e_{i}\rangle\, e_{i},
    \]
    and $e_{j}$ is included (i.e., $\Delta_{g_{j}} > 0$) if and only if $\|w_{j}\| > 0$.

    The cumulative span satisfies
    \begin{equation}
      \label{eq:Sigma-GS}
      \Sigma^{(c)}_{(b_{1},b_{2})}(n)
      \;=\;
      \dim \mathrm{span}\bigl(\mathcal{O}_{n}\bigr)
      \;=\;
      \mathrm{rank}\,
      \bigl[\rho_{c}(g)\,v_{0}\bigr]_{g \in B_{n}},
    \end{equation}
    where the right-hand side is the rank of the $q \times |B_{n}|$ matrix whose
    columns are the orbit vectors.
    This rank is independent of the enumeration order but depends on $v_{0}$.

    \begin{lemma}[Seed dependence of the rank]
      \label{lem:rank-seed}
      Fix a prime $q$ and $c \in (\mathbb{Z}/q\mathbb{Z})^{\times}$.
      \begin{enumerate}
        \item[(i)] For fixed~$n$, the rank
        $\Sigma^{(c)}_{(b_{1},b_{2})}(n) = \mathrm{rank}[\rho_{c}(g)\,v_{0}]_{g \in B_{n}}$
        is lower-semicontinuous in $v_{0} \in \mathbb{C}^{q} \setminus \{0\}$; it equals its
        maximal value $r_{\max}(c,n) \le \min(|B_{n}|, q)$ outside a proper algebraic subset of
        seeds.
        \item[(ii)] For every $n$ at least the diameter of $\mathrm{Cay}(G_q, S_q)$ and every
        $v_{0} \neq 0$, $\Sigma^{(c)}_{(b_{1},b_{2})}(n) = q$.
      \end{enumerate}
    \end{lemma}

    \begin{proof}
      (i) The rank is at least $k$ exactly when some $k \times k$ minor of the
      $q \times |B_n|$ orbit matrix is non-zero; the minors are polynomials in the entries
      of $v_0$, so $\{v_0 : \mathrm{rank} \ge k\}$ is open and, for $k = r_{\max}(c,n)$,
      the complement of the common zero set of those minors, a proper algebraic subset.
      The bound $\min(|B_n|, q)$ need not be attained: for instance the orbit vectors of
      $g$ and $gz$, with $z$ central, are proportional.
      (ii) For $n$ at least the diameter, $B_n = G_q$.
      The span of $\{\rho_c(g)\,v_0 : g \in G_q\}$ is a non-zero $\rho_c$-invariant subspace of
      $\mathbb{C}^q$, hence all of $\mathbb{C}^q$, since the Weil representation $\rho_c$ with
      non-trivial central character is irreducible~\cite{Beau2026a20}.
    \end{proof}

    \begin{remark}[Former Corollary~2.5 withdrawn]
      \label{cor:D-asymptotic}
      Versions up to~1.2 stated, as Corollary~2.5, that for generic $v_0$
      $\Sigma^{(c)}_{(b_{1},b_{2})}(n) \sim A(c, b_{1}, b_{2})\, n^{\dc}$ asymptotically, with
      $D(c, b_{1}, b_{2}, n) \to A(c, b_{1}, b_{2}) / A_{\mathrm{ref}}(c)$ absorbed into $\rcq$.
      This is withdrawn: $\Sigma^{(c)}$ is bounded by $q$ and equals $q$ from the diameter on
      (Lemma~\ref{lem:rank-seed}), so it has no asymptotic power law, and $\dc$ is the decay
      exponent of $\sigma_c$ on a pre-saturation window, not a growth exponent of $\Sigma^{(c)}$.
    \end{remark}

    \noindent
    \emph{Consequence for the residual factor.}
    By Lemma~\ref{lem:rank-seed}(ii), for every block instance $(b_1,b_2)$ and every seed,
    $\Sigma^{(c)}_{(b_1,b_2)}(n) = \Sigma^{(q-c)}_{(b_1,b_2)}(n) = q$ for $n$ at least the
    diameter, so the limit~\eqref{eq:r-def} exists and $\rcq = 1$; likewise $D \to 1$ for any
    reference span that also reaches~$q$.
    For $\rcq$ as defined in Definition~\ref{def:D}, Cases~A--C of
    Section~\ref{sec:classification} therefore cannot occur: the fourth case holds, Open
    Problems~\ref{op:P2} and~\ref{op:P3} have a negative answer, and one takes $\sigmacan := \sigmapair$.
    A non-trivial residual factor could only be defined as a ratio on a pre-saturation window,
    where $\Sigma^{(c)}$ does depend on the seed (Lemma~\ref{lem:rank-seed}(i)); this paper does
    not define such a ratio, and Tests~P2 and~P3 remain meaningful only for it.

  \subsection{Three layers of dependence}

    We separate the dependence of $D(c, b_{1}, b_{2}, n)$ into three distinct
    contributions:
    \begin{enumerate}
      \item \emph{Basis dependence.}
      The choice of initial seed vector $v_{0}$ and of the Gram--Schmidt
      orthonormalisation order implicitly defines a reference frame in
      $\mathbb{C}^{q}$.
      Different frames produce different increments $\Delta\Sigma^{(c)}$ at small
      $n$; by Lemma~\ref{lem:rank-seed}(ii) these effects disappear once the span saturates at~$q$.
      \item \emph{Block-instance dependence.}
      The pair $(b_{1}, b_{2})$ specifies a concrete realisation of the block at
      character~$c$.
      Different instances may encounter the geometry of $\rho_{c}$ differently,
      producing instance-dependent amplitudes.
      \item \emph{Convention dependence.}
      The normalisation convention (e.g., dividing by shell size $|S_{n}|$ or by
      block dimension $q$) introduces a multiplicative factor that is independent
      of both basis and instance.
    \end{enumerate}

    \noindent
    These three layers are logically independent and must not be conflated.
    The remainder of this paper addresses only layers 1 and 2; convention
    dependence is fixed by a canonical choice in Section~\ref{sec:canonical}.

  \section{Metaplectic Phase Structure and the Origin of $r(c,q)$}
  \label{sec:phase}

  \subsection{Central-phase coherence in Weil representations}

    The Weil representation $\rho_{c}$ of $\Heis$ involves metaplectic phases
    arising from the central coordinate $\gamma$ of the group.
    O12~\cite{Beau2026a16} introduced the coherence length
    \[
      \ellgam(n)
      \;=\;
      \left|
        \frac{1}{|S_{n}|}\sum_{g \in S_{n}} e^{2\pi i c\,\gamma_{g}/q}
      \right|,
    \]
    where $\gamma_{g}$ is the central coordinate of $g \in S_{n}$.
    When $\ellgam(n) \approx 1$ (high coherence), the central phase acts as a
    nearly constant multiplicative factor; when $\ellgam(n) \ll 1$ (low
    coherence), it randomises path contributions.

  \subsection{Phase as origin of $r(c,q)$}

    \begin{proposition}[Informal]
      \label{prop:phase-heuristic}
      A non-trivial pre-saturation analogue of $\rcq$ would be the residue of a phase-alignment
      mismatch between the initial basis $(b_{1}, b_{2})$ and the natural
      metaplectic phase structure of $\rho_{c}$.
      In the high-coherence regime ($\ellgam \approx 1$), this mismatch manifests
      as a constant multiplicative factor on $\Sigma^{(c)}(n)$, which would produce
      integer values such as the reported $\rcq \in \{1, 2, 3, 4\}$ if those values were
      confirmed.
    \end{proposition}

    \begin{remark}
      Proposition~\ref{prop:phase-heuristic} is an informal structural claim about a pre-saturation
      factor; for the asymptotic factor of Definition~\ref{def:D} the residue is identically~$1$
      (Lemma~\ref{lem:rank-seed}), and the proposition does not apply to it.
      Its precise formulation would depend on the outcome of the tests in
      Sections~\ref{sec:testP2}--\ref{sec:testP3} for a pre-saturation factor.
      In Case~A (pure basis artefact), the claim is that the mismatch is entirely
      absorbed by the choice of basis; in Cases~B and~C, a component of the
      mismatch is intrinsic to the block instance or the representation geometry.
    \end{remark}

  \subsection{Connection to a canonical choice of basis}
    \label{subsec:meta-basis}

    A natural candidate for a canonical basis choice in Case~A is the
    \emph{metaplectic-adapted basis}: the eigenbasis of $\rho_{c}(\gamma)$,
    which diagonalises the central action and thus fixes the phase reference
    unambiguously.
    In this basis, the phase $e^{2\pi i c\,\gamma/q}$ contributes a diagonal
    factor to every group element, and no interference between different
    $\gamma$-sectors can arise.
    Case~A would be confirmed if $\rcq$ varied under P2, was stable under P3, and equalled $1$ in this basis for all
    primes tested; $\rcq = 1$ alone confirms nothing, since it also holds in the fourth case
    below.

  \section{Test P2: Dependence on the Initial Basis}
  \label{sec:testP2}

  \subsection{Protocol}

    Fix $(q, c, b_{1}, b_{2})$.
    We consider unitary changes of basis of $\mathbb{C}^{q}$ preserving the standard inner
    product: the seed vector $v_{0}$ is replaced by $U v_{0}$ for $U \in \mathrm{U}(q)$,
    and the Gram--Schmidt orthonormalisation order is varied independently.
    Record $\rcq$ for each choice.

    \begin{openproblem}[P2]
      \label{op:P2}
      For a pre-saturation residual factor, does it vary under changes of initial basis at fixed
      $(q, c, b_{1}, b_{2})$?
      (For the factor of Definition~\ref{def:D}, the answer is negative: $\rcq \equiv 1$.)
    \end{openproblem}

  \subsection{Expected outcomes for a pre-saturation factor}

    If $\rcq$ varies under P2: the factor has a basis-dependent component,
    consistent with Proposition~\ref{prop:phase-heuristic}.
    If $\rcq$ is stable under P2: the factor is insensitive to the basis choice
    and its origin must lie elsewhere (block-instance geometry or representation
    structure).

  \section{Test P3: Dependence on the Block Instance}
  \label{sec:testP3}

  \subsection{Protocol}

    Fix $(q, c)$ and impose the metaplectic-adapted basis as a rigid convention.
    Vary the block instance $(b_{1}, b_{2})$.
    Record $\rcq$ for each instance.

    \begin{openproblem}[P3]
      \label{op:P3}
      For a pre-saturation residual factor, does it vary under changes of block instance
      $(b_{1}, b_{2})$ at fixed $(q, c)$ and a fixed canonical basis choice?
      (For the factor of Definition~\ref{def:D}, the answer is negative: $\rcq \equiv 1$.)
    \end{openproblem}

  \subsection{Expected outcomes for a pre-saturation factor}

    If $\rcq$ varies under P3: the factor has a component intrinsic to the
    block-instance geometry, i.e., to how the concrete realisation $(b_{1}, b_{2})$
    encounters the structure of $\rho_{c}$.
    This would make $\rcq$ a discrete invariant of the pair
    $(b_{1}, b_{2}, \rho_{c})$ rather than of $(c, q)$ alone.

  \section{Classification of the Residual Factor}
  \label{sec:classification}


  For a pre-saturation factor, the combined outcomes of Tests P2 and P3 would lead to one of the three
  cases below.
  A fourth logical possibility---$\rcq$ stable under both transformations,
  hence an invariant of $(c, q)$ alone, including the trivial value $\rcq \equiv 1$---is
  not excluded, and it is the case that holds for Definition~\ref{def:D}: within its toy model, the $q = 29$ check of O17~\cite{Beau2026a21} verifies the
  raw-count block-independence identity; it computes no normalised $\rcq$, so it gives no evidence
  of a variation of $\rcq$.
  For $\rcq$ as defined in Definition~\ref{def:D}, this fourth case is the one that holds
  (Lemma~\ref{lem:rank-seed} and the consequence stated after Remark~\ref{cor:D-asymptotic}), and
  no canonicalisation is needed: one takes $\sigmacan := \sigmapair$.
  Independently, on O16's contract the conjugate ratio of the exact block observable is also~$1$.
  Tests P2 and P3, and Cases~A--C below, are meaningful only for a residual factor defined on a
  pre-saturation window, which this paper does not define; they concern the normalisation
  parametrised by $(b_{1}, b_{2})$, whose relation to the pipeline's block index $(c_1,c_2,c_3)$ is
  not given (O17).

  \begin{description}
    \item[Case A (pure basis artefact).]
    $\rcq$ varies under P2 and is stable under P3.
    The residual factor is entirely attributable to the initial basis choice;
    it carries no information about the block-instance geometry.
    The canonical observable is obtained by selecting the metaplectic-adapted
    basis.
    \item[Case B (discrete instance invariant).]
    $\rcq$ is stable under P2 and varies under P3.
    The residual factor is a discrete invariant of the block instance at fixed
    $(q, c)$; it reflects a structural property of how $\rho_{c}$ interacts with
    the concrete pair $(b_{1}, b_{2})$.
    The canonical observable is obtained by quotienting by this invariant.
    \item[Case C (mixed).]
    $\rcq$ varies under both P2 and P3.
    The residual factor admits a multiplicative decomposition
    \begin{equation}
      \label{eq:r-decomp}
      \rcq \;=\; r_{\mathrm{base}}(c, U, q)\;\cdot\; r_{\mathrm{inst}}(c, b_{1}, b_{2}, q),
    \end{equation}
    where $r_{\mathrm{base}}$ depends on the unitary frame $U$ and vanishes under the
    metaplectic-adapted choice $U = U^{\mathrm{meta}}$, and $r_{\mathrm{inst}}$ is a
    discrete invariant of the block instance at fixed $(q, c)$.
    The canonical observable is obtained by first eliminating $r_{\mathrm{base}}$ via
    a canonical basis choice, then quotienting by $r_{\mathrm{inst}}$.
  \end{description}

  \begin{remark}
    Case~C is the most general of the three; it is not supported by any reported
    measurement.
    Within its toy model, O17 shows only that any variation of $\rcq$ would reside in the
    normalisation $D$; it does not show that $\rcq$ varies, nor which component of a variation
    would be absorbed by a change of basis.
  \end{remark}

  \section{Construction of the Canonical Observable}
  \label{sec:canonical}

  \subsection{Canonical normalisation}

    For each case of the classification of Section~\ref{sec:classification}, which concerns a
    pre-saturation factor only (for Definition~\ref{def:D} the fourth case holds and one takes
    $\sigmacan := \sigmapair$), we define the reference span $\Sigma^{(c)}_{\mathrm{ref}}(n)$ in
    Definition~\ref{def:D} as follows.

    In \textbf{Case~A}, the reference span is computed in the metaplectic-adapted
    basis $U^{\mathrm{meta}}$, defined as the eigenbasis of $\rho_{c}(\gamma_{0})$
    for a fixed central element $\gamma_{0} \neq 0$ of $G_{q}$.
    This diagonalises the central-phase action and fixes the phase reference
    unambiguously.
    By construction, $D(c, v_{0}^{\mathrm{meta}}, n) / D(q-c, v_{0}^{\mathrm{meta}}, n) = 1$
    for all $n$ sufficiently large, so $\rcq = 1$.

    In \textbf{Case~B}, the basis is already stable under P2, and the
    instance invariant $r_{\mathrm{inst}}(c, b_{1}, b_{2}, q)$ would be a pre-saturation analogue of
    the ratio~\eqref{eq:r-def} evaluated at any fixed canonical basis choice (the asymptotic ratio
    itself equals~$1$).
    The reference span is the minimal span over all instances:
    \[
      \Sigma^{(c)}_{\mathrm{ref}}(n)
      \;=\;
      \min_{(b_{1},b_{2})} \Sigma^{(c)}_{(b_{1},b_{2})}(n).
    \]
    The canonical pair observable is $\sigmacan(n) = \sigmapair(n) / r_{\mathrm{inst}}(c,q)$.

    In \textbf{Case~C}, the decomposition~\eqref{eq:r-decomp} is applied in two steps:
    (i) project onto the metaplectic-adapted basis to remove $r_{\mathrm{base}}$;
    (ii) quotient by the residual $r_{\mathrm{inst}}$ identified via Test~P3.

  \subsection{Canonical pair observable}

    \begin{definition}[Canonical pair observable]
      \label{def:sigmacan}
      The \emph{canonical pair observable} is
      \[
        \sigmacan(n)
        \;=\;
        \sigma_{c}^{\mathrm{can}}(n)\;\sigma_{q-c}^{\mathrm{can}}(n),
      \]
      where $\sigma_{c}^{\mathrm{can}}(n)$ is the span observable computed with
      canonical normalisation as defined above.
    \end{definition}

    \begin{openproblem}[Pipeline independence]
      \label{prop:independence}
      On a pre-saturation window, is $\sigmacan(n)$ independent of the choice of initial basis and
      block instance, up to corrections that are $o(n^{-\dpair})$?
    \end{openproblem}

    \begin{remark}
      Versions up to~1.2 stated this independence as Proposition~7.2 with a proof sketch that
      depended on the outcome of Tests~P2 and~P3.
      The statement is not proved and is withdrawn as a result, kept as an open problem with its
      number.
      It is restated here on a pre-saturation window, since past saturation the observable vanishes
      and an asymptotic statement would be empty.
      For the factor of Definition~\ref{def:D}, Lemma~\ref{lem:rank-seed} gives $\rcq \equiv 1$ and
      one takes $\sigmacan := \sigmapair$; on the pre-saturation window the seed dependence of
      $\Sigma^{(c)}$ (Lemma~\ref{lem:rank-seed}(i)) is not controlled, and the question remains open.
      In Case~A it would follow from the definition of the metaplectic-adapted basis; in
      Cases~B and~C it would require the shell-independence of $r_{\mathrm{inst}}$.
    \end{remark}

  \section{Discussion}
  \label{sec:discussion}

  \subsection{Interpretation of $r(c,q)$ in terms of phase alignment}

    The classification of Section~\ref{sec:classification} gives a structural
    interpretation of $\rcq$ consistent with the metaplectic phase analysis
    of O12~\cite{Beau2026a16}.
    In all three cases, the integer-valuedness of the reported values $\rcq \in \{1,2,3,4\}$,
    if those values were confirmed, would reflect a counting of coherent phase contributions rather than a continuous
    phase value: it would enumerate the number of $\gamma$-sectors of $\rho_{c}$
    that are in phase with the initial basis or block instance.

    This interpretation is consistent with the O12 observation that the
    coherence length $\ellgam(n)$ governs the transition between regimes of
    near-constant central phase and randomised contributions.
    In the language of O12, high-coherence shells ($\ellgam \approx 1$) would produce
    the integer multiplicity that would become a non-trivial pre-saturation factor;
    low-coherence shells would produce none.

    In the decomposition of Case~C, the component $r_{\mathrm{base}}$ admits a
    natural interpretation in terms of phase misalignment between the seed frame and
    the metaplectic structure of $\rho_{c}$, while $r_{\mathrm{inst}}$ admits a
    natural interpretation in terms of instance multiplicity, i.e., how many
    $\gamma$-sectors the concrete block $(b_{1}, b_{2})$ activates coherently.
    This decomposition \emph{suggests} these interpretations without requiring them:
    the classification would be determined empirically, and the interpretation would follow
    from the outcome of Tests~P2 and~P3 for a pre-saturation factor, which are open.

  \subsection{Implications for O20 and O21}

    The canonical observable $\sigmacan(n)$ proposed here would be the input
    required for the two subsequent papers; since a non-trivial canonicalisation would concern a
    pre-saturation factor not defined here, and pipeline independence in the pre-saturation window
    is open (Open Problem~\ref{prop:independence}), those papers cannot import it as a constructed,
    convention-free observable.
    O20 will investigate a persistence criterion and its comparison with the
    historical comparison range $\dpair \in [7.4, 10.6]$; the amplitude normalisation is irrelevant for the
    exponent but is needed to ensure that the persistence criterion is evaluated
    on a convention-free observable.
    O21 will examine the historical $\dpair \to \bstar$ prescription at the fibre level;
    Definition~\ref{def:sigmacan} gives the proposed starting point; for the factor of
    Definition~\ref{def:D} it reduces to $\sigmapair$.

  \subsection{What O19 does not resolve}

    Three directions remain open after this paper.

    \begin{openproblem}
      \label{op:arithmetic}
      Whether a non-trivial $\rcq$ occurs for any normalisation of interest, and, if so,
      the arithmetic origin of the reported values $\rcq \in \{1,2,3,4\}$---possibly
      related to quadratic Gauss sums and character multiplicities of the Weil
      representation---are left open.
    \end{openproblem}

    \begin{openproblem}
      \label{op:general-q}
      The stability of the canonical normalisation as $q \to \infty$ is not
      addressed; establishing it would require control of $\ellgam(n)$ in the
      large-$q$ limit.
    \end{openproblem}

    \begin{openproblem}
      \label{op:beyond-heis}
      The generalisation of the canonical normalisation to groups beyond
      $\Heis$ is deferred.
    \end{openproblem}

  \section{Stability of the Scaling Exponent}
  \label{sec:stability}


  \begin{proposition}[Exponent invariance]
    \label{prop:exponent}
    The capacity exponent $\dpair$ is unchanged under all canonicalisations of
    Section~\ref{sec:canonical}.
  \end{proposition}

  \begin{proof}
    The canonical normalisation modifies $\sigmapair(n)$ by a multiplicative
    factor that is either constant (Cases~A and~B) or decomposes into constant
    factors (Case~C).
    A constant multiplicative factor shifts $\log \sigmapair(n)$ by an additive
    constant and therefore does not affect the slope
    $\dpair = -\frac{\mathrm{d}\log\sigmapair(n)}{\mathrm{d}\log n}$ on the pre-saturation window
    (past saturation $\sigmapair$ vanishes and no slope is defined).
  \end{proof}

  \begin{remark}
    This proposition confirms that O19 operates strictly at the amplitude level.
    Statements of O16--O18 about the pair exponent and its comparison with the historical range
    $\delta \in [7.4, 10.6]$ are unaffected by a constant normalisation factor; the historical
    value near $7.44$ is not asserted by O16.
  \end{remark}

  \section{Conclusion}
  \label{sec:conclusion}


  We have made explicit the implicit pipeline normalisation factor
  $D(c, b_{1}, b_{2}, n)$ of the O12/O13 Gram--Schmidt span tracker and
  formulated two independence tests for the residual amplitude factor $\rcq$.
  For the factor as defined, an asymptotic limit, Lemma~\ref{lem:rank-seed} gives $\rcq \equiv 1$
  for every seed and block instance, so the fourth case holds and one takes $\sigmacan := \sigmapair$; a
  non-trivial factor would need a pre-saturation definition, for which the tests are open.
  The former Corollary~2.5, an asymptotic power law for $\Sigma^{(c)}$, is withdrawn.

  The classification (Cases~A, B, or~C) would determine the appropriate
  construction for the canonical pair observable $\sigmacan(n)$; its independence
  of pipeline conventions is an open problem (Open Problem~\ref{prop:independence}), not a
  result.
  A constant multiplicative factor leaves the capacity exponent $\dpair$ unchanged
  (Proposition~\ref{prop:exponent}).
  On the real O12/O13 pipeline, O16's conjugate-block identity forces $\rcq \equiv 1$ within
  its contract; the reported values $\rcq \in \{1,2,3,4\}$ remain a historical hypothesis.

  The normalisation step of the fibre-level programme therefore remains open.
  O20 and O21 refer to the observable proposed here, not to a constructed one.

  \clearpage
  \bibliographystyle{plain}
      \bibliography{cosmochrony-bibliography}

\end{document}
